\documentclass{standalone}
\usepackage{tikz}
\usetikzlibrary{intersections, calc}
\makeatletter\AtBeginDocument{\@buildstepset{1,2,3,4,5,6}}\makeatother
\usepackage{hsbcommon}

\begin{document}

% Define each curve's function exactly once, so \path and \draw
% can both call it by name instead of repeating the formula.
% x runs 1900 -> 2100 mapped to 0 -> 8 (25 years per unit).
% Heights are unitless index values, as in the original World3
% "standard run" chart (Meadows et al., Limits to Growth, 1972) --
% only the qualitative shape/ordering matters, not the numbers.
\pgfmathdeclarefunction{resources}{1}{\pgfmathparse{1.4 + 6.6/(1+exp(1.3*(#1-2.1)))}}      % nonrenewable resources: sigmoid decline
\pgfmathdeclarefunction{foodpc}{1}{\pgfmathparse{0.9 + 3.4*exp(-((#1-3.1)^2)/3.2)}}        % food per capita: early, moderate hump
\pgfmathdeclarefunction{indout}{1}{\pgfmathparse{0.45 + 3.6*exp(-((#1-3.4)^2)/2.0)}}       % industrial output per capita: peaks just after food
\pgfmathdeclarefunction{population}{1}{\pgfmathparse{6.4*exp(-((#1-4.6)^2)/10.0)}}         % population: tallest, broadest, latest peak, slow tail
\pgfmathdeclarefunction{pollution}{1}{\pgfmathparse{0.05 + 3.0*exp(-((#1-4.0)^2)/1.0)}}    % pollution: sharp, narrow, lags industry

\begin{tikzpicture}[every plot/.style={smooth},every node/.style={font=\sffamily}]

        \draw (0,0) -- (8.2,0);
        \draw (0,0.1) -- (0,-0.1) node[below,rotate=90,anchor=east] {1900};
        \draw (8,0.1) -- (8,-0.1) node[below,rotate=90,anchor=east] {2100};

%% define all the paths

        \path[domain=0:8, name path=resourcespath] plot ({\x},{resources(\x)});
        \path[domain=0:8, name path=foodpath]       plot ({\x},{foodpc(\x)});
        \path[domain=0:8, name path=indpath]        plot ({\x},{indout(\x)});
        \path[domain=0:8, name path=poppath]        plot ({\x},{population(\x)});
        \path[domain=0:8, name path=pollpath]       plot ({\x},{pollution(\x)});

        \path [name intersections={of=poppath and resourcespath,by=E}];

%% draw steps

\ifstep{1}{
    \begin{scope}[red]
        \draw[domain=0:8] plot ({\x},{resources(\x)}) node[above left, anchor=east,fill=white,font=\footnotesize] {resources};
    \end{scope}
}

\ifstep{2}{
    \begin{scope}[orange!90!black]
        \draw[domain=0:8] plot ({\x},{indout(\x)}) node[above left, anchor=east,fill=white,font=\footnotesize] {industrial output/capita};
    \end{scope}
}

\ifstep{3}{
    \begin{scope}[green!55!black]
        \draw[domain=0:8] plot ({\x},{foodpc(\x)}) node[above left, anchor=east,fill=white,font=\footnotesize] {food per capita};
    \end{scope}
}


\ifstep{4}{
    \begin{scope}[black]
        \draw[domain=0:8] plot ({\x},{population(\x)}) node[above left, anchor=east,fill=white,font=\footnotesize] {population};
    \end{scope}
}

\ifstep{5}{
    \begin{scope}[brown!85!black]
        \draw[domain=0:8] plot ({\x},{pollution(\x)}) node[above left, anchor=east,fill=white,font=\footnotesize] {pollution};
    \end{scope}
}

\ifstep{6}{
    \begin{scope}[blue]
        % \filldraw (E) circle (2pt);
        \draw[dashed] (E|-,0) node[below] {overshoot} -- (E|-,7);
    \end{scope}
}

    \end{tikzpicture}
\end{document}